% !TEX root = ../metatime_monograph.tex
\chapter{CKM CP Violation}
\label{ch:ckm_cp_violation}

\section{Direct Geometric Phase}

The CKM CP-violating phase \(\delta\) is not a derived quantity from the
Jarlskog invariant but a direct geometric angle of the tetrahedron:

\begin{equation}
\delta_{\text{CKM}} = \arccos\frac{L_4}{L_5}
                    = \arccos\frac{2}{5}
                    = 66.42^\circ.
\label{eq:ckm_delta}
\end{equation}

This angle is the dihedral angle between the annihilation plane
(dimension \(L_4 = 2\)) and the intention axis (dimension \(L_5 = 5\))
in the NOEMA tetrahedron.

\section{Comparison with PDG}

The PDG value is \(\delta_{\text{CKM}}^{\text{PDG}} = 65.6^\circ \pm 1.5^\circ\).

\[
\sigma = \frac{|66.42^\circ - 65.6^\circ|}{1.5^\circ} = 0.55.
\]

The prediction is within \(1\sigma\) of the experimental measurement.

\section{Jarlskog Invariant Consistency Check}

The Jarlskog invariant for the CKM matrix satisfies:

\[
J_{CP} = |V_{us}||V_{cb}||V_{ub}||V_{cs}|\sin\delta.
\]

Using our values:

\[
\begin{aligned}
\sin\delta &= \sin(66.42^\circ) = 0.9165, \\
J_{CP} &= 0.22485 \times 0.04082 \times 0.003628
          \times 0.97357 \times 0.9165 \\
       &= 2.97 \times 10^{-5}.
\end{aligned}
\]

The structural formula for \(J_{CP}\) (Eq.~\eqref{eq:jcp} in
Chapter~\ref{ch:ckm_matrix}) gives \(J_{CP} = 3.11 \times 10^{-5}\)
(PDG \(3.08 \times 10^{-5}\)).
The small discrepancy of \(4.5\%\) between the direct product and the
structural formula suggests a minor higher-order correction to the
CKM magnitudes, but the angle \(\delta = \arccos(L_4/L_5)\) itself is
robust.

\section{Why the Earlier Derivation Failed}

The earlier approach computed \(\delta\) by inverting the Jarlskog
relation:

\[
\sin\delta = \frac{J_{CP}}{|V_{us}||V_{cb}||V_{ub}||V_{cs}|}.
\]

This gave \(\delta = 73.6^\circ\) because the denominator was too small
(the \(|V_{cb}|\) formula gives \(0.04082\) while the PDG value is
\(0.04182\)).  The direct geometric formula \(\delta = \arccos(L_4/L_5)\)
bypasses this sensitivity and captures the CP phase as a fundamental
tetrahedron angle.
